\documentclass[12pt]{jlreq}
\usepackage{amsmath, amssymb}

\newcommand{\C}{\mathbb{C}}
\newcommand{\R}{\mathbb{R}}
\newcommand{\Z}{\mathbb{Z}}
\newcommand{\N}{\mathbb{N}}
\newcommand{\F}{\mathbb{F}}
\newcommand{\Prob}{\mathbb{P}}
\newcommand{\Exp}{\mathbb{E}}
\newcommand{\dd}{\,d}
\newcommand{\Ker}{\operatorname{Ker}}
\newcommand{\Impart}{\operatorname{Im}}
\newcommand{\Hom}{\operatorname{Hom}}

\begin{document}

\(2\) 次元単位球面
\[
  S^2=\{(x,y,z)\in\R^3\mid x^2+y^2+z^2=1\}
\]
上のいたるところでゼロでない \(C^\infty\) 級 \(2\)-形式
\[
  \omega=x\,dy\wedge dz+y\,dz\wedge dx+z\,dx\wedge dy
\]
と，\(S^2\) 上の \(C^\infty\) 級関数
\[
  H(x,y,z)=x^2+2y^2+3z^2
\]
を考える。\(S^2\) の \(C^\infty\) 級ベクトル場 \(X\) を
\[
  (dH)(Y)=\omega(X,Y),\qquad Y\in TS^2,
\]
で定義し，それが生成する \(1\)-径数微分同相群を \(\phi_t\) とする。

\begin{enumerate}
  \item 関数 \(H\) が \(\phi_t\) で不変であることを示せ。
  \item \(S^2\) の点を \(\phi_t\) によって次の \(3\) 種類に分類する。
  \begin{enumerate}
    \item 停留点，すなわち，すべての \(t\in\R\) に対し \(\phi_t(P)=P\) となる点 \(P\)。
    \item 周期点，すなわち，停留点ではないが，ある \(T>0\) に対し \(\phi_T(P)=P\) をみたす点 \(P\)。
    \item 停留点でも周期点でもない点。
  \end{enumerate}
  球面 \(S^2\) のおのおのの点が，(a), (b), (c) のいずれであるかを理由とともに答えよ。
  \item 停留点でも周期点でもない点 \(P\) に対し，\(\displaystyle\lim_{t\to\infty}\phi_t(P)\) および \(\displaystyle\lim_{t\to-\infty}\phi_t(P)\) を決定せよ。
\end{enumerate}

\end{document}